\chapter{Process Discovery Simulation Methodology}
\label{chp:process_discovery_methodology}

\textit{erklärt, wie aus den Daten ein ausführbares Model wird und wie es arbeitet? | Modellkonstruktion und Ausführungsmechanismen erklären}

\section{Process Discovery}
Process Discovery describes the process of extracting the control flow model of a given event log. 
The discovery follows an algorithms, that learns the control flow in different ways, that can be grouped into a family like structure. Each Family follows a methodology of how from observed sequences, a model structure is constructed. The categories are overlapping in their calculation methods. 
\cite{DiscoveryBenchmark}.

PETRI NET
The result of a process discovery can be expressed in a petri net as shown in figure ...[TBC INSERT FIGURE OF BASIC PETRI NET] with a starting point, an endpoint and several transsitions and places namely activities, constructing the process backbone. 

ALGORITHM DIFFERENCE IN DOSCOVERY
A petri net displays the learned process behaviour, that can be different, depending on the used algorithm. 

DOMAIN RULES / MODEL ADAPTIONS/ REPAIR
As the control flow will serve as the base for the following simulation models learning and execution, it should display the real processes domain rules, that can be specified by domain experts before the discovery to guide the discovery process [SOURCE - Rules guided discovery] or based on the resulting petri net, a so called process repair [SOURCE Massi. 2025] 

    \textit{
    \begin{enumerate}
        \item Candidates and Configuration of the different Discovery Algorithms
        Inductive Miner / IMf

        Heuristics Miner
        
        ILP Miner
        
        Split Miner

        \item Procedures for Rule checks and ggf. Model adaptions
    \end{enumerate}}

\section{Prosits Simulation Methodology}
WIE LERNT UND AREBITET PROSIT?

ProSiT is a probablistic white box simulation model that is...
It takes a petri net and an event log as its input and calculates several models accumulating to be one overal simulation model, that can simulate its learned behaviour.
It uses Probablistic Decision trees, with distributions in its leaves, rather than static values, which gives it a closer to reality set of resulting values for durations and actual final values. 
The user can decide on th granularity of the learned trees, by setting the maximum tree depth and the minimum amount of values per leaf.

EXPLAIN THE DIFFERENT TIMEMODELS
ProSiT constructs X seperate models, to generate its simulations.
Servcetime model

Waiting time model 

Inter Arrival time Model

Ressource Calendar (Model)?


Prosit has several Functional-Hyperarameters that can be configures when starting the simulation discovery. 
Activiating the Multitaksing, enables Prosit is able to recognise resources and their capacities, by learning simultaniously executed activities at the same resource. 

Activating the Workload Features, a queing behaviour can be learned, that can have an affect on the affected time models. 

[ADD MORE HYPERPARAMETERS]

    \begin{enumerate}
        \item Ankünfte, Routing, Ressourcenzuordnung, Kalender und Kapazitäten
        \item gelernte Zeitmodelle und Workload-Features
        \item Reihenfolge der Entscheidungen während der Simulation
        \item zusätzliche Diagnoseausgaben und Verifikationsprüfungen
    \end{enumerate}

\textit{Hier gehört ein nachvollziehbares Ausführungsbeispiel hinein: Ein Truck wird aktiviert, benötigt einen bestimmten Stop, erhält eine Ressource und durchläuft die einzelnen Zeitkomponenten. Jede Komponente muss zur Implementierung passen. So beantwortest du die Kritik an der bisher zu abstrakten technischen Darstellung.}